%------------------------------------------------------------------------------%
\section{Teste da Raiz}

O Teste da Raiz é similar ao Teste da Razão, ele é útil quando calcular
a raiz $n$-ésima for mais fácil do que calcular a razão ente os termos da série.

%------------------------------------------------------------------------------%
\begin{teorema}{Teste da Raiz}{teo:teste-raiz}
  Seja\index{Teste!raiz}
  \(\;
    \sum_{n=1}^\infty a_n
  \;\)
  uma série tal que $a_n\geq 0$ e
  \[
    \lim_{n\to \infty} \sqrt[n]{a_n} = \rho
  \]
  então,
  \begin{enumerate}
    \item se $\rho<1$ a série converge;
    \item se $\rho>1$ ou $\rho=\infty$ a série diverge;
    \item se $\rho=1$ o teste é inconclusivo.
  \end{enumerate}
\end{teorema}
%------------------------------------------------------------------------------%
\begin{proof}
  %----------------------------------------------------------------------------%
  \emph{Parte 1}

  Nesse caso, temos que $\rho<1$, assim podemos escolher $\varepsilon>0$
  tal que $\rho+\varepsilon<1$.
  Como
  \[
    \lim_{n\to \infty} \sqrt[n]{a_n}=\rho
  \]
  existe $N\in \N$ tal que se $n\geq N$
  \[
      \abs{\sqrt[n]{a_n}-\rho\,}\;<\;\varepsilon
      \quad\Rightarrow\quad
      \sqrt[n]{a_n} \;<\; \varepsilon+\rho
      \quad\Rightarrow\quad
      a_n \;<\; (\varepsilon+\rho)^n      \,.
  \]
  Dessa forma, pelo teste da comparação, temos que a série
  \[
    \sum_{n=1}^\infty a_n
  \]
  converge, pois seu termo geral, $a_n$ está limitado pelo termo geral da série geométrica
  \[
    \sum_{n=1}^\infty (\varepsilon+\rho)^n
  \]
  que é convergente pois $\varepsilon+\rho<1$.

  %----------------------------------------------------------------------------%
%  \emph{Parte 2}
%
%  {\color{red} Escrever}

\end{proof}
%------------------------------------------------------------------------------%

Assim como no Teste da Razão, no Teste da Raiz o teorema é inconclusivo
quando $\rho=1$. Observe as séries
\[
  \sum_{n=1}^\infty \frac{1}{n}
  \qquad \text{ e } \qquad
  \sum_{n=1}^\infty \frac{1}{n^2}
\]
em ambos os casos temos que
\(
  \lim \sqrt[n]{a_n}=1
\)
entretanto, a série harmônica diverge e a série~$p$, com $p=2$, converge.
%Fica como exercício verificar que
%\[
%  \lim \sqrt[n]{\frac{1}{n}\;} = \lim \sqrt[n]{\frac{1}{n^2}\;} = 1
%\]

Note o Teste da Raiz continua válido mesmo trocando a hipótese de que
\[
  \lim \sqrt[n]{a_n} = \rho
\]
por
\begin{equation}
  \label{eq:teste-raiz-variacao}
  \lim \sqrt[n+k]{a_n} = \rho \qquad \forall k \in \N     \,.
\end{equation}
Para ver que isso é verdade basta notar na demonstração que o termo geral
da série continua limitado pelo termo geral de uma série geométrica.

Seguem agora alguns exemplos do uso desse teste.

%------------------------------------------------------------------------------%
\begin{example}
  Use o teste da raiz para analisar a convergência da série
  \[
    \sum_{n=1}^\infty \frac{4^n}{\,(3n)^n\,}         \,.
  \]
  \solution
  Seja
  \[
    a_n = \frac{4^n}{\,(3n)^n\,}
  \]
  então
  \[
    \sqrt[n]{a_n\,} = \sqrt[n]{\frac{4^n}{\,(3n)^n\,}}
                    = \frac{4}{3n}
                    \to 0 < 1
  \]
  portanto a série converge.
\end{example}
%------------------------------------------------------------------------------%

%------------------------------------------------------------------------------%
\begin{example}
  Use o teste da raiz para analisar a convergência da série
  \[
    \sum_{n=1}^\infty\frac{1}{n^{n+1}}     \,.
  \]
  \solution
  Observe que, fazendo uma mudança de índice, podemos reescrever a série como
  \[
    \sum_{n=1}^\infty\frac{1}{\, n^{n+1}\,} =
    \sum_{n=2}^\infty\frac{1}{\,(n-1)^{n}\,}       \,.
  \]
  Podemos agora tomar
  \[
    a_n = \frac{1}{\,(n-1)^{n}\,}
  \]
  que nos leva a
  \[
    \sqrt[n]{a_n} = \frac{1}{n-1} \to 0 < 1
  \]
  o que nos permite concluir que série converge.
\end{example}
%------------------------------------------------------------------------------%

%Outra forma de resolver o exemplo anterior é usar a variação do
%Teste da Raiz~\eqref{eq:teste-raiz-variacao}.
%Neste caso, podemos tomar
%\[
%  a_n = \frac{1}{n^{n+1}}
%\]
%que nos leva a
%\[
%  \sqrt[n+1]{a_n} = \frac{1}{n} \to 0 < 1   \,.
%\]
%Assim, chegamos a conclusão de que a série converge.

%------------------------------------------------------------------------------%
\begin{example}
  Use o teste da raiz para analisar a convergência da série
  \[
    \sum_{n=1}^\infty \left[\ln\left(e^2+\frac{1}{n}\right)\right]^{n+1}    \,.
  \]
  \solution
  Fazendo uma transformação no índice podemos reescreve a série da seguinte forma
  \[
    \sum_{n=1}^\infty \left[\ln\left(e^2+\frac{1}{n}  \right)\right]^{n+1} \;=\;
    \sum_{n=2}^\infty \left[\ln\left(e^2+\frac{1}{n-1}\right)\right]^{n}   \,.
  \]
  Dessa forma temos que
  \[
    a_n=\left[\ln\left(e^2+\frac{1}{n-1}\right)\right]^{n}   \,.
  \]
  Aplicando o Teste da Raiz temos
  \[
    \sqrt[n]{a_n} = \ln\left(e^2+\frac{1}{n-1}\right)
                  \;\to\; \ln \left(e^2\right)
                  = 2 > 1      \,.
  \]
  Portanto essa série diverge.
\end{example}
%------------------------------------------------------------------------------%

%------------------------------------------------------------------------------%
